\documentclass{article}

\usepackage[utf8]{inputenc}     % pdfLaTeX
\usepackage[T1]{fontenc}
\usepackage[magyar]{babel}

\usepackage{amsmath}
\usepackage{amsthm}
\usepackage{amssymb}
\usepackage{tikz}
\usepackage{siunitx}   % mértékegységek, hőmérséklet
\usetikzlibrary{arrows.meta, positioning}

\theoremstyle{definition}
\newtheorem{definition}{Definíció}

\title{Hőtan}
\author{Rovenszky Robin, Simon László}
\date{Legutóbb frissítve: 2026. Március 12.}

\begin{document}

\maketitle
\tableofcontents
\newpage

\section{Alap definíciók}
\begin{definition}[Hőállapot]
    Részecskék hőmozgása.
\end{definition}

\begin{definition}[Hőérzet]
    Szubjektív, befolyásolható
\end{definition}

\subsection{Hőmérséklet mérése}

\subsubsection{Fahrenheit}
A Fahrenheit egy tapasztalati hőmérsékleti skála. A 100F a feltalálójának a testhőmérséklete volt.
\subsubsection{Celsius}
A Celsius is egy tapasztalati hőmérsékleti skála.
%Drawing missing
\begin{itemize}
    \item \SI{100}{\degreeCelsius} -- a víz forráspontja
    \item \SI{0}{\degreeCelsius} -- a víz fagyáspontja
\end{itemize}
\subsubsection{Kelvin}
A Kelvin abszolút hőmérsékleti skála, tehát a 0K az abszolút nulla hőmérséklet.
\subsubsection{Hőmérők}
%Drawing missing
\begin{itemize}
    \item $212^\circ F$ = \SI{100}{\degreeCelsius} $\approx 373^\circ K$
    \item $32^\circ F$ = \SI{0}{\degreeCelsius} $\approx 273^\circ K$
    \item \SI{1}{\degreeCelsius} = $(^\circ F-32) \times \frac{5}{9}$
\end{itemize}

\subsection{Hőtágulás}

%Kísérlet, ábra missing
\begin{itemize}
    \item Sín \begin{itemize}
        \item Melegben locsolják
        \item Közök hagyása - zakatolás
    \end{itemize}
    \item Villanyvezetékek - Betonsúlyok
    \item Csövek %insert tube picture here
    \item Hidak - Görgők %insert bridge picture here
\end{itemize}

Mitől függ?
\begin{equation}
    \Delta l = \alpha \cdot l_{0} \cdot \Delta T \\
\end{equation}
$\Delta l$ - hosszváltozás \\
$l_{0}$ - eredeti hossz \\
$\Delta T$ - hőmérsékletváltozás \\
$\alpha$ - anyagi minőség\\ \\
Hőtágulási együttható:
\begin{equation}
    [\alpha]=\frac{1}{^{\circ}C} = \frac{1}{K}
\end{equation}


%Note to self: draw function drawn in class with tikz

\section{Párolgás vagy forrás?}
Mindkettő folyadék $\rightarrow$ gáz\\
Mi a különbség?\\
Párolgás: csak a felületen, bármilyen hőmérsékleten\\
Forrás: jól meghatározott hőmérsékleten, "forráspont". A folyadék belselyéből is.
\\
Mitől függ a párolgás? 
\begin{itemize}
    \item anyagi minőségtől
    \item folyadék felszínének méretétől
    \item hőmérséklettől
    \item páratartalom
    \item szél, légnyomás
\end{itemize}

Hőt von el a környezetétől.
\[
Q = L_p \cdot m
\]

\section{Lecsapódás}
gáz $\rightarrow$ folyékony\\
Hidegebb felület

\end{document}